\documentclass[a4paper]{article}
\usepackage[utf8]{inputenc}
\usepackage[russian,english]{babel}
\usepackage[T2A]{fontenc}
\usepackage[left=10mm, top=20mm, right=18mm, bottom=15mm, footskip=10mm]{geometry}
\usepackage{indentfirst}
\usepackage{amsmath,amssymb}
\usepackage{mathastext} 
\usepackage[italicdiff]{physics}
\usepackage{graphicx}
\usepackage{caption}
\usepackage{float}
\usepackage{caption}
\renewcommand{\thefootnote}{\fnsymbol{footnote}}
\usepackage{tablefootnote}
\usepackage{footmisc}
\usepackage{textcomp}
\usepackage{multicol}
\usepackage{multirow}
\usepackage[parfill]{parskip}
\usepackage{makecell}
\usepackage{array}  
\usepackage{siunitx}
\usepackage{booktabs}
\sisetup{
    exponent-product = \cdot,
    per-mode = symbol,
    separate-uncertainty = true,
    uncertainty-separator = \,
}
\renewcommand\cellalign{cc} 
\renewcommand\cellgape{\Gape[2pt]} 
\usepackage[utf8]{inputenc}\newcommand{\approxtext}[1]{\ensuremath{\stackrel{\text{#1}}{\approx}}}
\usepackage[colorlinks=true, urlcolor=blue, linkcolor=blue, citecolor=blue]{hyperref}
\graphicspath{{images/}}
\DeclareGraphicsExtensions{.pdf,.png,.jpg}
\usepackage{wrapfig}
\captionsetup{labelformat=empty}
\usepackage{caption}
\captionsetup[figure]{name=Рисунок}
\captionsetup[table]{name=Таблица}
%\usepackage{newtxtext,newtxmath} 
\title{\textbf{Отчет о выполнении лабораторной работы 4.7.3} \\[0.5em] \large\itshape Изучение поляризованного света}
\date{}
\author{Котляров Михаил, Б01-402}

\begin{document}

\maketitle
	
\section{Введение}

\subsubsection*{Цель работы:} ознакомление с методами получения и анализа
поляризованного света.

\subsubsection*{Оборудование и инструментальные погрешности}

оптическая скамья с осветителем; зелёный
светофильтр; два поляроида; чёрное зеркало; полированная эбонитовая пластинка;
стопа стеклянных пластинок; слюдяные пластинки разной толщины; пластинки в 1/4
и 1/2 длины волны; пластинка в одну длину волны для зелёного света (пластинка
чувствительного оттенка).


\section*{Описание экспериментальной установки}
Экспериментальная установка собрана на оптической скамье и включает в себя: источник белого света с конденсорной системой, зелёный светофильтр, два поляроида с градуированными лимбами (поляризатор и анализатор), чёрное зеркало на поворотном столике, полированную эбонитовую пластинку, стопу стеклянных пластинок, набор двоякопреломляющих слюдяных пластинок различной толщины, пластинки с фазовым сдвигом $\lambda/4$ и $\lambda/2$, а также пластинку чувствительного оттенка (пластинка $\lambda$ для зелёной спектральной компоненты). Все оптические элементы закреплены на рейтерах, позволяющих перемещать их вдоль скамьи и вращать вокруг оптической оси. Наблюдение интенсивности и цвета прошедшего света осуществляется визуально или с помощью фотоприёмника, расположенного в конце оптической оси.

\section*{Ход работы и результаты измерений}

\subsection*{1. Определение разрешённых направлений поляроидов}
Установим на оптической скамье осветитель, первый поляроид и чёрное зеркало. Будем вращать поляроид вокруг направления луча, а чёрное зеркало вокруг вертикальной оси методом последовательных приближений, добиваясь минимальной яркости отражённого пятна. Затем заменим зеркало на второй поляроид и, скрестив поляроиды, определим его разрешённое направление.
\textbf{Результаты:} Получили положения запрещённых направлений: для первого поляроида вертикальное направление соответствует углу $88 \pm 1^\circ$, для второго поляроида -- $86 \pm 1^\circ$. Погрешность учтена только систематическая прибора.

\subsection*{2. Измерение угла Брюстера и показателя преломления эбонита}
Установим на скамью вместо чёрного зеркала полированную эбонитовую пластинку. Будем измерять угол Брюстера по лимбу, фиксируя положение минимальной интенсивности отражённого луча сначала без светофильтра, затем с зелёным фильтром. Рассчитаем показатель преломления по формуле $n = \tan \theta_{\text{Б}}$. Оценим случайную и систематическую погрешности измерений.
\textbf{Результаты:} Получили следующие значения удвоенного угла Брюстера: $2\theta_{\text{Б}} = 127^\circ, 130^\circ, 128^\circ$ (без фильтра) и $2\theta_{\text{Б}} = 126^\circ, 129^\circ, 127^\circ$ (с фильтром). Исходя из 6 измерений, случайная погрешность равна $0{,}601^\circ$, а с учётом систематической погрешности ($1^\circ$) итоговая погрешность составила $1{,}17^\circ$. Среднее значение угла Брюстера составило $\langle \theta_{\text{Б}} \rangle = 63,92^\circ \pm 1{,}17^\circ$.  Показатель преломления эбонита определили как $n = \tan \langle \theta_{\text{Б}} \rangle = 2{,}04 \pm 0,11$. Табличный показатель преломления приблизительно 1,50–1,66.

\subsection*{3. Исследование поляризации света в лучах от стопы Столетова}
Установим вместо эбонитовой пластинки стопу стеклянных пластинок, наклонив её под углом Брюстера к падающему пучку. Осветим стопу неполяризованным светом. Будем анализировать отражённый и преломлённый лучи с помощью вращаемого поляроида, фиксируя углы максимальной и минимальной интенсивности.
\textbf{Результаты:} Определили, что вектор $\vec{E}$ в отражённом луче поляризован под углом $97^\circ$ от вертикали.

\subsection*{4. Определение главных направлений двоякопреломляющих пластинок}
Поместим исследуемые двоякопреломляющие пластинки поочерёдно между скрещенными поляроидами. Будем вращать каждую пластинку вокруг направления луча, наблюдая за интенсивностью прошедшего света, и зафиксируем углы, при которых интенсивность минимальна. При повороте у обеих пластин по четыре раза за
оборот достигался минимум интенсивности прошедшего света. Он соотвествовал
совпадению главных осей пластин с разрешенными направлениями поляроидов.
\textbf{Результаты:} Минимум интенсивности для первой пластинки наблюдали при $28^\circ$, для второй -- при $13^\circ$. Эти положения соответствуют совпадению главных направлений пластинок с разрешёнными направлениями поляроидов.

\subsection*{5. Выделение пластинок $\lambda/4$ и $\lambda/2$}
Добавим в схему зелёный светофильтр. Установим разрешённое направление первого поляроида горизонтально, а главные направления исследуемых пластинок -- под углом $45^\circ$ к горизонтали. С помощью второго поляроида проанализируем поляризацию света на выходе из каждой пластинки, вращая анализатор и фиксируя изменения интенсивности.
\textbf{Результаты:} При установке второй пластинки на угол $58^\circ$ ($13^\circ + 45^\circ$) интенсивность при вращении анализатора практически не менялась, что указало на круговую поляризацию; следовательно, вторая пластинка является пластинкой $\lambda/4$. При установке первой пластинки на угол $73^\circ$ ($28^\circ + 45^\circ$) вместо ожидаемого минимума при скрещенных поляроидах наблюдали максимум, что свидетельствует об инверсии плоскости поляризации; таким образом, первая пластинка имеет толщину $\lambda/2$.

\subsection*{6. Определение «быстрой» и «медленной» осей в пластинке $\lambda/4$}
Поместим между скрещенными поляроидами пластинку чувствительного оттенка и убедимся в её пурпурной окраске при освещении белым светом. Добавим к схеме пластинку $\lambda/4$, совмещая её главные направления с направлениями пластинки чувствительного оттенка. Будем поворачивать систему на $180^\circ$ вокруг вертикальной оси и наблюдать изменение цвета прошедшего света.

Мозаичная пластинка состоит из участков разной толщины $d_i$, поэтому каждый квадрат имеет свою оптическую разность хода $\Gamma_i = d_i |n_o - n_e|$ и, соответственно, свой спектральный коэффициент пропускания:
\begin{equation}
	I_i(\lambda) = I_0 \sin^2(2\alpha)\sin^2(\frac{\pi \Gamma_i}{\lambda}),
\end{equation}
где $\alpha$ – угол между разрешённым направлением поляризатора и главной осью данного участка.

\textbf{Результаты:} При совпадении «быстрых» осей обеих пластинок разность хода увеличилась до $5\lambda/4$, что привело к погашению красной части спектра и изменению окраски на зеленовато-голубую. При перпендикулярном расположении быстрых осей разность хода составила $3\lambda/4$, погасилась фиолетово-голубая часть спектра, и цвет стал оранжево-жёлтым. По характеру изменения цвета определили ориентацию быстрой и медленной осей в исследуемой пластинке $\lambda/4$.

\subsection*{7. Интерференция поляризованных лучей}
Расположим между скрещенными поляроидами мозаичную слюдяную пластинку, состоящую из участков разной оптической толщины. Будем последовательно вращать саму мозаичную пластинку и второй поляроид, наблюдая за изменением окраски и интенсивности в отдельных квадратиках.
\textbf{Результаты:} При вращении пластинки каждый квадратик сохранял свой цвет, однако глобально наблюдались разные оттенки для разных участков. При вращении анализатора квадратики приобретали различные цвета, что связано с изменением проекций когерентных волн на направление пропускания анализатора и зависимостью интерференционной картины от разности фаз, создаваемой участками разной толщины.

\subsection*{8. Определение направления вращения светового вектора в эллиптически поляризованной волне}
Установим разрешённое направление первого поляроида под углом $10^\circ$--$20^\circ$ к горизонтали, а второго -- вертикально. Вращая пластинку $\lambda/4$, добьёмся минимальной интенсивности на выходе, получив эллиптически поляризованный свет с вертикально ориентированной малой осью. Для определения направления вращения поместим дополнительную пластинку $\lambda/4$ с горизонтально ориентированной быстрой осью. Будем анализировать ориентацию результирующей плоскости поляризации на выходе из системы.
\textbf{Результаты:} После прохождения через систему пластинок световой вектор перешёл в смежные квадранты, что указало на однонаправленное вращение эллипсов поляризации в обеих пластинках. На основе известного направления вращения для эталонной пластинки $\lambda/4$ определили, что полученная эллиптическая волна обладает правой круговой поляризацией.

\section{Выводы}

В ходе выполнения лабораторной работы были изучены основные методы получения и анализа поляризованного света, а также исследованы явления, связанные с поляризацией при отражении, преломлении и прохождении через анизотропные среды.

\subsection*{Основные результаты}
\begin{enumerate}
    \item \textbf{Определение разрешённых направлений поляроидов.} Методом последовательных приближений с использованием чёрного зеркала установлены запрещённые направления: $88 \pm 1^\circ$ и $86 \pm 1^\circ$ для первого и второго поляроидов соответственно. Результаты согласуются между собой в пределах погрешности, что подтверждает корректность методики.
    
    \item \textbf{Измерение показателя преломления эбонита.} Среднее значение угла Брюстера составило $\langle \theta_{\text{Б}} \rangle = 63{,}92^\circ \pm 1{,}17^\circ$, что дало показатель преломления $n = \tan \langle \theta_{\text{Б}} \rangle = 2{,}04 \pm 0{,}11$. Полученное значение существенно превышает табличный диапазон для эбонита ($1{,}50$--$1{,}66$). Вероятные причины расхождения:
    \begin{itemize}
        \item субъективность визуального определения минимума интенсивности отражённого луча;
        \item возможное отклонение поверхности эбонитовой пластинки от идеальной плоскости или наличие микрошероховатостей;
        \item влияние рассеянного света и неидеальной коллимации пучка;
        \item возможное окисление или загрязнение поверхности, изменяющее оптические свойства границы.
    \end{itemize}
    Тем не менее, методика демонстрирует принципиальную возможность определения $n$ для непрозрачных диэлектриков.
   
    
    \item \textbf{Идентификация фазовых пластинок.} Пластинки $\lambda/4$ и $\lambda/2$ успешно выделены по характеру изменения интенсивности при вращении анализатора: для $\lambda/4$ интенсивность практически постоянна (круговая поляризация), для $\lambda/2$ наблюдаются чёткие минимумы и максимумы с инверсией плоскости поляризации. Результаты соответствуют теоретическим предсказаниям для фазовых сдвигов $\pi/2$ и $\pi$.
    
    \item \textbf{Определение быстрой и медленной осей.} С помощью пластинки чувствительного оттенка установлено, что при совпадении быстрых осей $\lambda$- и $\lambda/4$-пластинок наблюдается зеленовато-голубой цвет (погашение красной компоненты, $\Gamma = 5\lambda/4$), а при перпендикулярной ориентации -- оранжево-жёлтый (погашение сине-фиолетовой, $\Gamma = 3\lambda/4$). Качественное согласие с теорией интерференции поляризованных лучей подтверждено.
    
    \item \textbf{Интерференция в мозаичной пластинке.} При вращении самой пластинки цвет отдельных квадратиков не меняется (множитель $\sin^2(2\alpha)$ не зависит от $\lambda$), а при вращении анализатора наблюдается инверсия цветов вследствие фазового сдвига на $\pi$. Наблюдения полностью соответствуют формуле интенсивности $I \propto \sin^2(\pi\Gamma/\lambda)$ для скрещенных поляроидов.
    
    \item \textbf{Направление вращения вектора $\vec{E}$.} Методом компенсации фаз с использованием эталонной $\lambda/4$-пластинки определено, что эллиптически поляризованная волна обладает правой круговой поляризацией. Результат получен в полном соответствии с алгоритмом, описанным в методических указаниях.
\end{enumerate}

\subsection*{Оценка точности и источники погрешностей}
Основной вклад в неопределённость результатов вносят:
\begin{itemize}
    \item \textbf{Систематическая погрешность отсчёта по лимбу} ($\pm 1^\circ$) -- ограничивает точность угловых измерений во всех пунктах работы;
    \item \textbf{Субъективность визуальной регистрации минимумов интенсивности} -- особенно критична при определении угла Брюстера и главных направлений пластинок;
    \item \textbf{Неидеальная монохроматичность источника} -- даже с зелёным фильтром спектр имеет конечную ширину, что размывает интерференционные эффекты для фазовых пластинок;
    \item \textbf{Неточная юстировка оптической схемы} -- малые угловые рассогласования элементов приводят к систематическим сдвигам в измерениях;
    \item \textbf{Качество поверхностей оптических элементов} -- шероховатости, загрязнения, неоднородности толщины пластинок влияют на воспроизводимость результатов.
\end{itemize}

Случайная погрешность, оценённая по серии измерений угла Брюстера, составила $0{,}601^\circ$, что свидетельствует о хорошей воспроизводимости методики при многократных измерениях. Однако доминирующая систематическая составляющая ($1^\circ$) ограничивает итоговую точность.

\subsection*{Заключение}
Цель работы достигнута: освоены практические методы получения линейно, эллиптически и циркулярно поляризованного света, а также методы анализа состояния поляризации с использованием поляроидов и фазовых пластинок. Полученные результаты подтверждают основные положения теории поляризации света, а выявленные погрешности носят методический характер и могут быть уменьшены при использовании более точной регистрирующей аппаратуры и улучшении качества оптических поверхностей.

\end{document}